% ============================================================
%  CAPITOLO 2 --- ANALISI DEI REQUISITI
% ============================================================
\chapter{Analisi dei requisiti}

\raggedbottom

\section{Intervista}

Di seguito si riporta il testo dell'intervista con il committente, che ha definito
i requisiti del sistema informativo. Le righe sono numerate per facilitare il
riferimento nelle sezioni successive.

\begin{tcolorbox}[colback=lightgray, colframe=primary!40, left=8pt, right=8pt,
                  top=6pt, bottom=6pt]
\setlength{\parskip}{6pt}

\textbf{1.} Il proprietario o i proprietari (di seguito: \textit{proprietario}) possono
effettuare l'accesso al gestionale tramite credenziali (nome, cognome, email, password).

\textbf{2.} Si vuole tenere traccia di clienti, prenotazioni, prodotti in magazzino,
ordini, fornitori, dipendenti e turni.

\textbf{3.} Per ogni \textit{cliente} nel gestionale, si memorizzano nome, cognome,
numero di telefono ed email.

\textbf{4.} I clienti possono prenotare \textit{ombrelloni} o \textit{lettini} oppure
possono sottoscrivere un \textit{abbonamento} valido per tutta la stagione.

\textbf{5.} Il prezzo dell'abbonamento stagionale è calcolato in base alla posizione
dell'ombrellone, mentre il prezzo della prenotazione è calcolato in base a posizione e durata.

\textbf{6.} Di ogni ombrellone è necessario memorizzare la posizione, il numero di lettini
e il numero di sedie associati. Il numero di lettini è standard: 2 per ogni ombrellone.

\textbf{7.} Le informazioni che legano ombrelloni e clienti (data inizio, data fine)
sono conservate nelle \textit{prenotazioni}.

\textbf{8.} Il proprietario può gestire i \textit{turni} dei dipendenti, visualizzare
la forza lavoro disponibile in determinate giornate e modificare i turni
per coprire l'intera giornata.

\textbf{9.} È possibile verificare la disponibilità in \textit{magazzino} ed effettuare
\textit{ordini} per singoli o gruppi di prodotti.

\textbf{10.} Il pagamento degli ordini viene effettuato alla consegna; una volta aggiunto
un ordine, il contenuto viene automaticamente aggiunto al magazzino.

\end{tcolorbox}

% -- 2.2 Ambiguità -------------------------------------------
\section{Rilevamento delle ambiguità e correzioni proposte}

A seguito della lettura dei requisiti sono state individuate le seguenti ambiguità
e si propongono le relative correzioni.

\begin{table}[H]
  \centering
  \caption{Tabella delle ambiguità rilevate e correzioni proposte}
  \label{tab:ambiguita}
  \rowcolors{2}{rowalt}{white}
  \begin{tabularx}{\textwidth}{C{1cm} L{3.5cm} L{3.5cm} L{5.5cm}}
    \toprule
    \rowcolor{primary!15}
    \textbf{Riga} & \textbf{Termine originale} & \textbf{Termine adottato} & \textbf{Motivazione} \\
    \midrule
    3 & lettino & \textit{Lettino} (entità) & Disambiguato da \textit{numeroLettini}
        (attributo di Ombrellone): il lettino è un'entità indipendente prenotabile
        separatamente dall'ombrellone. \\[4pt]
    3 & abbonamento & \textit{Abbonamento} (entità distinta da Prenotazione) &
        L'abbonamento differisce dalla prenotazione per la durata (stagionale)
        e per la struttura della tariffa. Si modella come entità separata. \\[4pt]
    4 & posizione & \textit{fila} + \textit{numero} (attributi composti) &
        La posizione di un ombrellone è identificata da una fila e un numero
        progressivo all'interno della fila. Si decompone in due attributi atomici. \\[4pt]
    6 & turno & \textit{Turno} (entità) con orario inizio e fine &
        Un turno è caratterizzato da data, orario di inizio e orario di fine.
        Un dipendente può avere più turni nella stessa giornata. \\[4pt]
    7 & prodotto & \textit{Prodotto} (entità in magazzino) distinto da
        \textit{Fornitore} & Aggiunto il concetto di \textit{Fornitore} come entità
        autonoma, poiché un prodotto può essere acquistato da fornitori diversi. \\[4pt]
    \bottomrule
  \end{tabularx}
\end{table}

% -- 2.3 Specifiche ristrutturate ----------------------------
\section{Specifiche in linguaggio naturale ristrutturato}

A seguito dell'analisi, si riporta il testo ristrutturato con i termini univoci
adottati. I concetti principali sono evidenziati in \textsc{maiuscoletto}.

\medskip
\begin{tcolorbox}[colback=lightgray, colframe=primary!40, left=8pt, right=8pt,
                  top=6pt, bottom=6pt]
\setlength{\parskip}{6pt}

Il \entita{Proprietario} accede al gestionale tramite credenziali (nome, cognome,
email, password). Nel sistema possono essere presenti più proprietari.

Il proprietario può registrare \entita{Clienti}, specificando nome, cognome,
numero di telefono ed email. Un cliente può essere associato a una o più
\entita{Prenotazioni} oppure a uno o più \entita{Abbonamenti}.

Una \entita{Prenotazione} lega un \entita{Cliente} a un \entita{Ombrellone}
e/o a un \entita{Lettino} per un periodo definito da \attributo{dataInizio}
e \attributo{dataFine}. Il prezzo viene calcolato in base alla durata e alla
posizione dell'ombrellone.

Un \entita{Abbonamento} è simile a una prenotazione ma ha durata stagionale
e una struttura tariffaria diversa (calcolata su base di durata e posizione).
Anche l'abbonamento mostra solo il prezzo senza gestire il pagamento.

Di ogni \entita{Ombrellone} si memorizzano: \attributo{fila}, \attributo{numero},
\attributo{numeroLettini}, \attributo{numeroSedie}.

Un \entita{Dipendente} è caratterizzato da nome, cognome, codice fiscale e
recapito. Il proprietario può registrare i \entita{Turni} dei dipendenti,
ciascuno con \attributo{data}, \attributo{orarioInizio} e \attributo{orarioFine}.

In \entita{Magazzino} vengono gestiti i \entita{Prodotti}, ciascuno con
\attributo{nome}, \attributo{quantità} e \attributo{prezzoUnitario}. Quando
la quantità scende sotto una soglia, il sistema segnala la mancanza.

Il proprietario può effettuare \entita{Ordini} verso i \entita{Fornitori},
specificando i prodotti e le quantità desiderate. Al momento della consegna,
le quantità ordinate vengono aggiunte automaticamente al magazzino.

\end{tcolorbox}

% -- 2.4 Concetti principali ---------------------------------
\section{Estrazione dei concetti principali}

\begin{table}[H]
  \centering
  \caption{Concetti principali estratti dall'analisi}
  \label{tab:concetti}
  \rowcolors{2}{rowalt}{white}
  \begin{tabularx}{\textwidth}{L{3cm} C{2cm} L{8.5cm}}
    \toprule
    \rowcolor{primary!15}
    \textbf{Termine} & \textbf{Tipo} & \textbf{Descrizione} \\
    \midrule
    \entita{Proprietario}    & Entità      & Utente del gestionale, accede tramite credenziali \\
    \entita{Cliente}         & Entità      & Persona registrata a cui si assegnano prenotazioni o abbonamenti \\
    \entita{Ombrellone}      & Entità      & Postazione identificata da fila e numero \\
    \entita{Prenotazione}    & Associazione & Lega un cliente a un ombrellone/lettino per un periodo \\
    \entita{Abbonamento}     & Entità      & Contratto stagionale cliente-ombrellone con tariffa dedicata \\
    \entita{Dipendente}      & Entità      & Lavoratore dello stabilimento \\
    \entita{Turno}           & Entità      & Fascia oraria lavorativa di un dipendente \\
    \entita{Prodotto}        & Entità      & Articolo presente in magazzino \\
    \entita{Fornitore}       & Entità      & Azienda da cui si acquistano i prodotti \\
    \entita{Ordine}          & Entità      & Richiesta di fornitura verso un fornitore \\
    \relazione{Contiene}     & Associazione & Lega un ordine ai prodotti ordinati (con quantità) \\
    \bottomrule
  \end{tabularx}
\end{table}
